\section{Programmation du mode F1 -- convoyeur/guichet seul}
\label{secModeF1Restreint}

\begin{UPSTIwarning}{Version restreinte du mode F1}
Le mode~F1 décrit en \ref{secModeF1} couple le comportement du
convoyeur/guichet, celui du \textbf{robot} et celui de l'\textbf{IHM}. Ici,
on ne programme \textbf{que la partie convoyeur/guichet} de ce mode : ni le
robot, ni la signalisation détaillée sur IHM ne sont traités dans ce
travail. Le guichet doit donc être commandé de façon autonome, sans
attendre une quelconque action du robot.
\end{UPSTIwarning}


\subsection{Niveau 1 -- Capacité unitaire du guichet}

\begin{UPSTIPreparation}[Cahier des charges -- niveau 1]
On souhaite, dans un premier temps, commander les seuls bloqueurs
(\texttt{BE}, \texttt{BP}) de façon à garantir la \textbf{capacité
unitaire} du guichet (\ref{secModeF1}), \textbf{sans distinguer} le cas
d'une palette pleine de celui d'une palette vide : toute palette entrée
dans le guichet y est immobilisée jusqu'à ce que le poste soit à nouveau
libéré manuellement par le montage suivant (aucune temporisation n'est
demandée à ce niveau). La colonne lumineuse du guichet
(\texttt{CR}/\texttt{CV}/\texttt{CB}) reste \textbf{éteinte}.

\UPSTIquestion{Etablir la commande permettant de garantir la capacité unitaire du guichet.}

\end{UPSTIPreparation}

% \subsection{Niveau 2 -- Maintien temporisé des palettes pleines}

% \begin{UPSTIPreparation}[Cahier des charges -- niveau 2]
% On complète le comportement précédent : une palette \textbf{portant un
% gobelet} (couple \texttt{aid}/\texttt{dgt} $=(1,1)$, cf.
% \ref{secDonneesGuichet}) doit désormais être \textbf{maintenue en poste
% pendant une durée \texttt{T\#3s}} avant d'être libérée, alors qu'une
% palette \textbf{vide} ($(1,0)$) continue d'être libérée sans délai. La
% capacité unitaire du guichet établie au niveau~1 doit rester garantie en
% toute circonstance.

% \UPSTIquestion{Modifier le grafcet du niveau~1 pour introduire la
% temporisation \texttt{T\#3s} sur les seules palettes portant un gobelet,
% sans remettre en cause la capacité unitaire du guichet.}

% \UPSTIquestion{Préciser la condition de branchement permettant de
% distinguer les deux cas (palette pleine / palette vide) à partir des
% capteurs \texttt{aid} et \texttt{dgt}.}
% \end{UPSTIPreparation}

% \begin{UPSTIManipulation}{Implémentation}
%     \UPSTIetape{Créer les sections \texttt{F1\_Guichet} (SFC principal) et
%     \texttt{F1\_Guichet\_Actions} (actions associées).}

%     \UPSTIetape{Implémenter le grafcet du niveau~2 (qui intègre le
%     niveau~1).}

%     \UPSTIetape{Rédiger, dans une section ST dédiée, la \textbf{synthèse
%     des sorties} \texttt{m\_nqxBE}, \texttt{m\_nqxBP},
%     \texttt{m\_nqxCR}, \texttt{m\_nqxCV} et \texttt{m\_nqxCB} à partir des
%     étapes actives du grafcet (la colonne lumineuse restant éteinte,
%     conformément au cahier des charges de ce travail).}

%     \UPSTIetape{Faire vérifier avant de tester sur l'installation.}
% \end{UPSTIManipulation}
